\documentclass[a4paper,11pt]{article}

% --- Packages ---
\usepackage[T1]{fontenc}
\usepackage{geometry}
\geometry{margin=1in}
\usepackage{amsmath, amssymb}
\usepackage{graphicx}
\usepackage[hidelinks]{hyperref}
\usepackage{booktabs}
\usepackage{xcolor}
\usepackage{tikz}
\usepackage{pgfplots}
\pgfplotsset{compat=1.18}
\usetikzlibrary{shapes.geometric, arrows.meta, positioning, shadows, fit, calc, backgrounds}
\usepackage{float}
\usepackage{enumitem}
\usepackage{fancyhdr}
\usepackage{listings}
\usepackage{setspace}
\usepackage{caption}
\usepackage{tcolorbox}
\usepackage{parskip}
\usepackage{algorithm}
\usepackage{algpseudocode}

% --- Style Settings ---
\hypersetup{
    colorlinks=true,
    linkcolor=black,
    filecolor=magenta,
    urlcolor=blue,
    pdftitle={GHOST Protocol Litepaper v1.0},
}

\definecolor{ghostblue}{RGB}{20, 40, 80}
\definecolor{codegray}{rgb}{0.5,0.5,0.5}
\definecolor{codepurple}{rgb}{0.58,0,0.82}
\definecolor{backcolour}{rgb}{0.95,0.95,0.92}

\lstdefinestyle{mystyle}{
    backgroundcolor=\color{backcolour},
    commentstyle=\color{green!60!black},
    keywordstyle=\color{magenta},
    numberstyle=\tiny\color{codegray},
    stringstyle=\color{codepurple},
    basicstyle=\ttfamily\footnotesize,
    breakatwhitespace=false,
    breaklines=true,
    captionpos=b,
    keepspaces=true,
    numbers=left,
    numbersep=5pt,
    showspaces=false,
    showstringspaces=false,
    showtabs=false,
    tabsize=2
}
\lstset{style=mystyle}

\pagestyle{fancy}
\fancyhf{}
\setlength{\headheight}{14pt}
\rhead{\textcolor{gray}{GHOST Protocol v1.0}}
\lhead{\textcolor{gray}{Litepaper}}
\cfoot{\thepage}

% --- Title Page ---
\title{
    \vspace*{3cm}
    \textbf{\huge GHOST Protocol} \\[0.5em]
    \Large Privacy-Preserving Rate Discovery for Decentralised Lending \\[0.2em]
    \large \textit{Litepaper v1.0}
    \vspace{2cm}
}
\author{Snehendu Roy \quad Debanjan Mondal \quad Abhirup Banerjee}
\date{2025}

\begin{document}

% --- Page 1: Dedicated Title Page ---
\maketitle
\thispagestyle{empty}

\begin{abstract}
    \noindent Decentralised lending today suffers from a fundamental tension: rate
    transparency enables market efficiency but invites strategic manipulation,
    front-running, and collusion. We present \textbf{GHOST} (\textbf{G}eneral
    \textbf{H}osting \textbf{o}f \textbf{S}ealed \textbf{T}icks), a protocol
    that resolves this tension by combining sealed-bid discriminatory-price
    auctions with confidential compute execution. Lenders submit rate bids
    encrypted under a secp256k1 public key held by a Chainlink Confidential
    Runtime Environment (CRE) node. Neither the storage layer, other
    participants, nor any single node operator can observe plaintext rates.
    Matching, rate validation, and fund disbursement occur entirely within the
    CRE's trusted execution boundary, producing a system where price discovery
    is both market-driven and manipulation-resistant. GHOST achieves what open
    order books cannot: \textit{truthful rate revelation without the attack
    surface of public bids.}
\end{abstract}

\clearpage
\setcounter{page}{2}

\tableofcontents
\vspace{2em}
\clearpage

\section{Introduction}

The dominant interest rate model in DeFi, pioneered by Aave and Compound,
computes borrowing rates algorithmically from pool utilisation:

\begin{equation}
    R_t = R_0 + R_{\text{slope}_1} \cdot \frac{U}{U_{\text{optimal}}}
    \qquad \text{if } U \leq U_{\text{optimal}}
    \label{eq:utilisation}
\end{equation}

This mechanism, while elegant, \textit{conflates liquidity risk with credit
risk}. All borrowers pay the same rate regardless of creditworthiness. The
rate reflects how much capital is available, not whether the borrower
deserves it. For protocols aspiring to support undercollateralised or
variable-collateral lending, where default probability is nonzero and
heterogeneous, this is insufficient.

Eli and Alexandre \cite{eli2025} proposed a tick-based auction framework
where lending pools are decomposed into discrete rate ticks, each
representing a sub-pool at a specific interest rate. Lenders select their
rate; borrowers consume ticks cheapest first. Their simulations demonstrate
convergence to a market-clearing rate that reflects the borrower's true
default probability, with speed proportional to information disclosure.
This framework builds on the classic credit rationing literature of
Stiglitz and Weiss \cite{stiglitz1981} and Bester \cite{bester1985}, which
established that asymmetric information between lenders and borrowers leads
to inefficient rate-setting under centralised intermediation.

However, their model assumes an \textit{open} auction: all bids are visible
on-chain throughout the book-building phase. While this enables competitive
rebidding, it also creates vulnerability to front-running, last-second
sniping, and coordinated rate manipulation, precisely the market frictions
the authors acknowledge as limitations.

\begin{tcolorbox}[colback=backcolour, colframe=gray, arc=2pt, boxrule=0.5pt]
\small\textbf{Core Contribution.}\quad
GHOST resolves this by making the auction \textbf{sealed}. Every rate bid,
both lender rates and borrower maximum rates, is encrypted client-side
under a public key whose corresponding private key exists \textit{only}
inside a Chainlink CRE workflow. The server that stores these bids is
architecturally incapable of reading them. Decryption, matching, and
settlement execute atomically within the CRE's confidential boundary,
inheriting the security guarantees of Chainlink's Decentralised Oracle
Network (DON) infrastructure.
\end{tcolorbox}

The result is a lending protocol with four guarantees:

\begin{itemize}[leftmargin=1.4em, itemsep=3pt]
  \item Rates are \textbf{market-discovered}, not algorithmically imposed.
  \item Each lender earns their own bid rate (\textbf{discriminatory
        pricing}), eliminating free-riding.
  \item \textbf{No participant}, including the protocol operator, observes
        rates before settlement.
  \item Fund custody is \textbf{non-custodial}, mediated by a
        privacy-preserving vault with compliance enforcement.
\end{itemize}

\section{Theoretical Foundation}

\subsection{Tick-Based Rate Decomposition}

We adopt the lending pool formalism of Eli and Alexandre \cite{eli2025}.
A lending pool $LP_{B}$ for borrower class $B$ consists of $n$ ticks
$\{r_1, r_2, \ldots, r_n\}$ with $0 < r_1 < r_2 < \cdots < r_n$, where
each tick $r_i$ holds deposit volume $d_i \geq 0$. The cumulative deposit
function is:

\begin{equation}
    CD_{r_i} = \sum_{j=1}^{i} d_j
    \label{eq:cumdeposit}
\end{equation}

A borrower seeking principal $K$ consumes ticks in ascending rate order.
The fill fractions $(\delta_1, \ldots, \delta_n)$ solve:

\begin{equation}
    \min_{\{\delta_i\}} \sum_{i} \delta_i \cdot r_i
    \qquad \text{subject to} \quad
    \sum_{i} \delta_i = K, \quad 0 \leq \delta_i \leq d_i
    \label{eq:optim}
\end{equation}

The borrower's effective blended rate is:

\begin{equation}
    R_{\text{eff}} = \frac{\sum_{i} \delta_i \cdot r_i}{K}
    \label{eq:blended}
\end{equation}

Under discriminatory pricing, each lender $k$ at tick $r_i$ receives
exactly $r_i$ on their matched principal, not the blended rate. This
preserves individual price signals and eliminates the payoff irrelevance
of infra-marginal bids that enables collusion in uniform-price auctions,
as documented by Myers et al.\ \cite{myers2021} and Goldreich
\cite{goldreich2007}.

\subsection{Lender Utility and Truthful Bidding}

Lender $k$ evaluates a private default probability $p_{k,B}$ for borrower
class $B$, with recovery rate $RR$. Their expected utility from depositing
$m$ tokens at tick $r_i$ is:

\begin{equation}
    u_{\text{lender},k}(m, r_i)
    = m \cdot \Bigl[(1 - p_{k,B}) \cdot r_i
                    - p_{k,B} \cdot (1 - RR)\Bigr]
      \cdot \mathbf{1}_{[CD_{r_i} \leq K]}
    \label{eq:lenderutil}
\end{equation}

The indicator captures the matching constraint: the lender earns only if
their tick is consumed. The lender's \textit{true-value rate} $r_k^*$, the
minimum rate yielding non-negative expected utility, satisfies:

\begin{equation}
    \boxed{r_k^* = \frac{p_{k,B} \cdot (1 - RR)}{1 - p_{k,B}}}
    \label{eq:truevalue}
\end{equation}

In an open auction with sufficient liquidity ($CD_{r^*+\epsilon} > K$),
lenders are incentivised to bid at $r_k^*$: underbidding risks negative
returns while overbidding risks exclusion. GHOST preserves this equilibrium
property while removing the information leakage inherent in open auctions.
This result extends the ascending-bid analysis of Ausubel \cite{ausubel2004}
and the competitive bidding theory of Milgrom and Weber \cite{milgrom1982}
to the sealed-bid lending context.

\subsection{The Sealed-Bid Extension}

In GHOST, lender bids are not observable during the book-building phase.
Each lender encrypts their rate $r_i$ as:

\begin{equation}
    c_i = \mathrm{ECIES\_encrypt}(pk_{\mathrm{CRE}},\; r_i)
    \label{eq:encrypt}
\end{equation}

where $pk_{\mathrm{CRE}}$ is the secp256k1 public key of the CRE workflow.
The storage layer observes only ciphertext $c_i$. The borrower similarly
encrypts their maximum acceptable rate:

\begin{equation}
    c_{\max} = \mathrm{ECIES\_encrypt}(pk_{\mathrm{CRE}},\; R_{\max})
    \label{eq:borrowerencrypt}
\end{equation}

The CRE decrypts both sides and validates $R_{\text{eff}} \leq R_{\max}$
within its confidential execution boundary.
\textbf{No party outside the CRE ever observes plaintext rates.}

In sealed-bid discriminatory auctions, truthful bidding remains a dominant
strategy when the number of bidders is sufficiently large relative to units
available \cite{monostori2014}. The discriminatory format further
discourages strategic bid shading: each lender's payoff depends solely on
their own bid, not on any clearing price. Klemperer \cite{klemperer2004}
provides a comprehensive treatment of the trade-offs between auction formats
that motivates this design choice.

\subsection{Borrower Utility and the Acceptance Mechanism}

The borrower's utility from accepting a proposal with blended rate
$R_{\text{eff}}$ is:

\begin{equation}
    u_{\text{borrower}} = K \cdot (ROI - R_{\text{eff}})
    \label{eq:borrowerutil}
\end{equation}

where $ROI$ is the borrower's expected return on borrowed capital. The
borrower accepts if $R_{\text{eff}} \leq R_{\max} \leq ROI$.

To prevent frivolous borrowing and ensure credible participation, GHOST
requires borrowers to post collateral $C \geq K \cdot \mu(\tau)$, where
$\mu(\tau)$ is a multiplier indexed by credit tier $\tau$:

\begin{table}[H]
\centering
\caption{Collateral multiplier schedule. Lower tiers require more
  collateral relative to principal, incentivising repayment and
  reputation-building.}
\renewcommand{\arraystretch}{1.4}
\begin{tabular}{@{}lcc@{}}
  \toprule
  \textbf{Credit Tier} & \textbf{Multiplier $\mu(\tau)$} & \textbf{Notes} \\
  \midrule
  Bronze   & $2.0\times$ & New borrowers \\
  Silver   & $1.8\times$ & Established history \\
  Gold     & $1.5\times$ & Strong repayment record \\
  Platinum & $1.2\times$ & Prime borrowers \\
  \bottomrule
\end{tabular}
\label{tab:tiers}
\end{table}

For cross-asset collateral (e.g.\ ETH against a USD-denominated loan),
the requirement adjusts via a real-time Chainlink price feed:

\begin{equation}
    C_{\mathrm{ETH}} \geq \frac{K \cdot \mu(\tau)}{P_{\mathrm{ETH/USD}}}
    \label{eq:crosscollateral}
\end{equation}

\section{Protocol Architecture}

GHOST implements a three-layer separation of concerns, each operating under
distinct trust assumptions.

\begin{table}[H]
\centering
\caption{Three-layer architecture overview.}
\renewcommand{\arraystretch}{1.4}
\begin{tabular}{@{}p{0.18\textwidth}p{0.22\textwidth}p{0.50\textwidth}@{}}
  \toprule
  \textbf{Layer} & \textbf{Component} & \textbf{Responsibilities} \\
  \midrule
  Layer 1: Custody & Chainlink CPT Vault &
    Shielded addresses, private off-chain transfers, withdrawal tickets,
    ACE compliance checks. GHOST never takes direct custody. \\
  Layer 2: Storage & GHOST API Server (Hono) &
    Stores encrypted blobs only, EIP-712 authentication, tracks loan
    lifecycle state, queues fund movements. Cannot decrypt rates or
    execute matching. \\
  Layer 3: Settlement & Chainlink CRE (DON) &
    Decrypts sealed rates, runs matching engine, disburses principal,
    monitors loan health. Holds $sk_{\mathrm{CRE}}$ as a
    threshold-encrypted DON secret. \\
  \bottomrule
\end{tabular}
\end{table}

\subsection{Layer 1: Privacy-Preserving Custody}

All fund custody is delegated to a Chainlink Compliant Private Transfer
vault on Ethereum Sepolia. This layer provides shielded addresses
(unlinkable recipient identifiers), private transfers validated off-chain
against the ACE PolicyEngine, withdrawal tickets redeemable on-chain within
a time-bound window, and compliance enforcement without exposing transaction
metadata publicly. The vault holds all deposited tokens; GHOST orchestrates
movements through a pool wallet that signs EIP-712 typed-data messages
against the vault's API.

The proof-of-concept operates with two synthetic assets registered on the
vault: \textbf{gUSD}, a USD-pegged stablecoin used as the primary lending
denomination, and \textbf{gETH}, a synthetic ether token used as borrower
collateral. Lenders deposit gUSD into tick positions at their chosen rate;
borrowers post gETH collateral valued against a live Chainlink ETH/USD
price feed. This two-token design captures the essential cross-asset
dynamics of real lending markets, including currency risk in collateral
valuation, liquidation threshold monitoring, and the need for
oracle-sourced pricing, while remaining tractable for confidential compute execution within
CRE's WASM runtime.

\subsection{Layer 2: Blind Storage}

The GHOST API server (Hono + Bun) functions as a deliberately blind storage
layer. It maintains in-memory state across deposit slots, lend intents,
borrow intents, match proposals, active loans, pending transfers, and credit
scores. All user-facing endpoints authenticate via EIP-712 typed-data
signatures with timestamp validation ($\pm 5$ minutes) to prevent replay
attacks. The server \textit{cannot} decrypt rates, execute matching, or
disburse funds autonomously.

\subsection{Layer 3: Confidential Settlement}

The Chainlink CRE serves as the protocol's settlement brain. Three
independent CRE workflows orchestrate the lending lifecycle:

\begin{table}[H]
\centering
\caption{CRE workflow schedule and responsibilities.}
\renewcommand{\arraystretch}{1.4}
\begin{tabular}{@{}p{0.22\textwidth}cc p{0.42\textwidth}@{}}
  \toprule
  \textbf{Workflow} & \textbf{Interval} & \textbf{Function} \\
  \midrule
  Rate Discovery and Matching & 30\,s epoch &
    Retrieves all pending intents, decrypts sealed rates using
    $sk_{\mathrm{CRE}}$, executes tick matching, emits proposals \\
  Transfer Execution & 15\,s cycle &
    Polls pending fund movements; executes EIP-712-signed private
    transfers through the vault API \\
  Loan Health Monitoring & 60\,s cycle &
    Checks active loans against collateralisation thresholds using
    real-time ETH/USD Chainlink price feeds \\
  \bottomrule
\end{tabular}
\label{tab:workflows}
\end{table}

\section{Matching Algorithm}

The epoch-based matching engine operates within CRE Workflow 1. Let
$\mathcal{L} = \{l_1, \ldots, l_m\}$ denote the active lend intents and
$\mathcal{B} = \{b_1, \ldots, b_q\}$ the pending borrow intents.

\begin{algorithm}[H]
\caption{Greedy Tick Matching (within CRE confidential boundary)}
\label{alg:matching}
\begin{algorithmic}[1]
  \State \textbf{Input:} $\mathcal{L}$ (lend intents, decrypted),
         $\mathcal{B}$ (borrow intents, decrypted)
  \State Sort $\mathcal{B}$ by $K_j$ \textsc{descending};
         sort $\mathcal{L}$ by $r_i$ \textsc{ascending}
  \For{each $b_j \in \mathcal{B}$}
    \State $\text{filled}_j \gets 0$;\;
           $\text{wSum}_j \gets 0$;\;
           $\mathcal{T}_j \gets \emptyset$
    \For{each $l_i \in \mathcal{L}$ with remaining capacity $a_i > 0$}
      \State $\mathrm{take} \gets \min(a_i,\; K_j - \text{filled}_j)$
      \State $\text{filled}_j \mathrel{+}= \mathrm{take}$;\quad
             $\text{wSum}_j \mathrel{+}= \mathrm{take} \cdot r_i$;\quad
             $a_i \mathrel{-}= \mathrm{take}$
      \State Append $(l_i,\, \mathrm{take},\, r_i)$ to $\mathcal{T}_j$
      \If{$\text{filled}_j = K_j$} \textbf{break} \EndIf
    \EndFor
    \State $R_{\mathrm{eff},j} \gets \text{wSum}_j \,/\, \text{filled}_j$
    \If{$\text{filled}_j = K_j$ \textbf{and}
        $R_{\mathrm{eff},j} \leq R_{\max,j}$}
      \State Emit proposal $P_j = (b_j,\; \mathcal{T}_j,\;
             R_{\mathrm{eff},j})$
    \Else
      \State Release all ticks in $\mathcal{T}_j$ back to $\mathcal{L}$
    \EndIf
  \EndFor
\end{algorithmic}
\end{algorithm}

The algorithm is greedy-optimal: sorting borrows largest-first maximises
liquidity utilisation, while sorting lends cheapest-first minimises
borrower cost. The blended rate $R_{\mathrm{eff}}$ is verified against the
borrower's encrypted maximum before any proposal is emitted.

\section{Incentive Mechanisms}

\subsection{Rejection Penalty}

A borrower who receives a valid proposal ($R_{\text{eff}} \leq R_{\max}$)
and rejects it incurs a penalty that compensates lenders for locked capital
during the proposal window:

\begin{equation}
    \text{penalty} = 0.05 \cdot C, \qquad \text{refund} = 0.95 \cdot C
    \label{eq:penalty}
\end{equation}

\subsection{Liquidation and Loss Distribution}

When a loan's health ratio breaches the liquidation threshold, the protocol
seizes collateral and distributes it pro-rata by principal contribution:

\begin{equation}
    \text{protocol fee} = 0.05 \cdot C, \qquad
    \text{lender}_i\text{ share}
    = 0.95 \cdot C \cdot \frac{\delta_i}{\displaystyle\sum_i \delta_i}
    \label{eq:liquidation}
\end{equation}

\section{Privacy Model}

GHOST achieves a layered privacy architecture where information is
compartmentalised by role. The critical invariant is:

\begin{tcolorbox}[colback=backcolour, colframe=gray, arc=2pt, boxrule=0.5pt]
\centering\small\bfseries
No single entity possesses both the encrypted rates and the decryption key.
\end{tcolorbox}

\begin{table}[H]
\centering
\caption{Information compartmentalisation across protocol layers.}
\renewcommand{\arraystretch}{1.4}
\begin{tabular}{@{}p{0.25\textwidth}ccc@{}}
  \toprule
  \textbf{Data} & \textbf{Public Chain} & \textbf{GHOST Server} & \textbf{CRE} \\
  \midrule
  Vault deposits        & Visible  & Hidden         & Hidden \\
  Transfer amounts      & Hidden   & Known          & Known \\
  Lender rates          & Hidden   & Encrypted blob & Plaintext at settlement \\
  Borrower max rate     & Hidden   & Encrypted blob & Plaintext at settlement \\
  Matched allocations   & Hidden   & Post-match     & Known at match time \\
  Loan terms            & Hidden   & Recorded       & Known \\
  \bottomrule
\end{tabular}
\label{tab:privacy}
\end{table}

\section{Role of Chainlink Infrastructure}

GHOST is designed as a native CRE application, leveraging multiple Chainlink
services as composable infrastructure primitives:

\begin{table}[H]
\centering
\caption{Chainlink services consumed by GHOST.}
\renewcommand{\arraystretch}{1.4}
\begin{tabular}{@{}p{0.28\textwidth}p{0.62\textwidth}@{}}
  \toprule
  \textbf{Chainlink Service} & \textbf{Role in GHOST} \\
  \midrule
  CRE Workflows &
    Confidential compute for rate decryption, matching, and settlement \\
  ConfidentialHTTPClient &
    End-to-end encryption for all CRE to external service communication \\
  Vault DON Secrets &
    Threshold-encrypted storage of $sk_{\mathrm{CRE}}$ and pool wallet
    key across DON nodes \\
  Chainlink Price Feeds &
    Real-time ETH/USD pricing for cross-asset collateral valuation \\
  Chainlink ACE &
    PolicyEngine compliance enforcement on every private transfer \\
  \bottomrule
\end{tabular}
\label{tab:chainlink}
\end{table}

\section{Discussion}

\subsection{Sealed vs.\ Open Auctions}

The Eli--Alexandre model achieves rate convergence through iterative
rebidding: lenders observe the order book and adjust rates competitively.
GHOST's sealed-bid design sacrifices this iterative convergence for
manipulation resistance. In practice, convergence emerges
\textit{across epochs} rather than within a single auction, as market
participants learn clearing rates from past loan originations and calibrate
future bids accordingly. This inter-epoch learning dynamic mirrors the
information aggregation properties documented in repeated sealed-bid
auctions \cite{milgrom1982}.

\subsection{Capital Efficiency}

The tick-based model requires total liquidity $CD_{r_n} > K$ for
competitive rate discovery. GHOST addresses this through continuous
matching: lend intents persist across epochs until matched or cancelled,
creating a standing liquidity book that grows over time. The credit tier
system further improves capital efficiency by reducing collateral
requirements for proven borrowers, unlocking leverage that would otherwise
remain idle.

\subsection{Limitations and Future Work}

The current implementation uses in-memory state, suitable for
proof-of-concept but requiring persistent storage for production
deployment. The credit scoring mechanism is endogenous and simple; future
iterations could integrate verifiable credential systems or zero-knowledge
proof-based credit attestations. Multi-token lending pools, variable
maturity structures, and secondary market trading of matched tick positions
represent natural extensions of the framework.

\section{Conclusion}

GHOST demonstrates that privacy and market efficiency in decentralised
lending are not mutually exclusive. By combining the tick-based rate
discovery framework of Eli and Alexandre \cite{eli2025} with Chainlink
CRE's confidential compute capabilities, the protocol achieves sealed-bid
price discovery with discriminatory settlement.

The three-layer architecture, comprising custody, blind storage, and
confidential settlement, ensures that no single component can unilaterally
observe rates, manipulate matching, or misappropriate funds. The sealed-bid
format resolves the front-running and collusion vulnerabilities inherent in
open order books while preserving the equilibrium properties of truthful
rate revelation established in the auction theory literature
\cite{milgrom1982, ausubel2004, klemperer2004}.

Future work will explore the integration of zero-knowledge proofs for
on-chain collateral verification, a lending-native vault with programmable
lock and release semantics, and extension of the credit-tier mechanism to
support fully undercollateralised lending for verified institutional
borrowers.

\begin{thebibliography}{99}
\small

\bibitem{eli2025}
C.\ Eli and H.\ Alexandre,
``Rate Discovery in Decentralised Lending,''
\textit{Journal of The British Blockchain Association},
vol.~8, no.~2, 2025.

\bibitem{stiglitz1981}
J.\ E.\ Stiglitz and A.\ Weiss,
``Credit Rationing in Markets with Imperfect Information,''
\textit{American Economic Review},
vol.~71, no.~3, pp.~393--410, 1981.

\bibitem{bester1985}
H.\ Bester,
``Screening vs.\ Rationing in Credit Markets with Imperfect Information,''
\textit{American Economic Review},
vol.~75, no.~4, pp.~850--855, 1985.

\bibitem{monostori2014}
Z.\ Monostori,
``Discriminatory versus Uniform Price Auctions,''
\textit{MNB Occasional Papers},
no.~111, 2014.

\bibitem{myers2021}
E.\ Myers, A.\ Bostian, and H.\ Fell,
``Asymmetric Cost Pass-Through in Multi-Unit Procurement Auctions:
An Experimental Approach,''
\textit{Journal of Industrial Economics},
vol.~69, pp.~109--130, 2021.

\bibitem{ausubel2004}
L.\ M.\ Ausubel,
``An Efficient Ascending-Bid Auction for Multiple Objects,''
\textit{American Economic Review},
vol.~94, no.~5, pp.~1452--1475, 2004.

\bibitem{milgrom1982}
P.\ Milgrom and R.\ J.\ Weber,
``A Theory of Auctions and Competitive Bidding,''
\textit{Econometrica},
vol.~50, pp.~1089--1122, 1982.

\bibitem{goldreich2007}
D.\ Goldreich,
``Underpricing in Discriminatory and Uniform Price Treasury Auctions,''
\textit{Journal of Financial and Quantitative Analysis},
vol.~42, no.~2, pp.~443--466, 2007.

\bibitem{klemperer2004}
P.\ Klemperer,
``Auctions: Theory and Practice,''
\textit{SSRN Electronic Journal}, 2004.

\end{thebibliography}

\end{document}
